\section{Introduction}
Atrial fibrillation is sustained on a remodelled conduction substrate:
interstitial fibrosis slows and disorganises propagation, and the resulting
wavelength shortening, source--sink mismatch, and unidirectional block permit
reentry. Because that substrate is a network of coupled excitable elements,
graph-theoretic and spectral summaries of the atrial conduction network have
been proposed as substrate biomarkers. Algebraic connectivity $\lam_2$, the
second-smallest Laplacian eigenvalue (Fiedler,
1973)~\cite{fiedler1973}, governs the stability of synchronised states in
coupled excitable and oscillator networks through the master-stability function
(Pecora and Carroll, 1998)~\cite{pecora1998} and the synchronizability
eigenratio $\lam_N/\lam_2$~\cite{barahona2002}, and has been applied as a
real-time seizure detector~\cite{bomela2020}. The sensitivity of $\lam_2$ to a
single conduction weight has a classical closed form,
$\partial\lam_2/\partial w_{ij}=(\phii_{2,i}-\phii_{2,j})^2$ (Ghosh and Boyd,
2006)~\cite{ghosh2006}, where $\phii_2$ is the Fiedler vector. Aggregating that
sensitivity over a fibrosis-weighted conduction-uncoupling field yields a
per-region Spectral Fragility Index (SFI). The identity is classical; the
contribution here is its evaluation.

What has not been established is whether such an index carries information
beyond the substrate features already in use. Fibrosis burden, fibrosis spatial
heterogeneity, fibrosis patch size, and connectivity measures such as the global
minimum cut, the percolation threshold, and $\lam_2$ itself are computed from
the same meshes at negligible cost. A spectral index is worth adopting only if
it improves discrimination when added to those features; only if the improvement
is large enough to matter, judged against a threshold fixed in advance rather
than a $p$ value alone; and only if it survives cross-validation that holds out
whole patients and shape families, so that anatomy cannot leak between training
and test folds. Evaluations meeting all three conditions have not been reported.

This study evaluated SFI as a strict add-on to six competitor features for the
prediction of simulator reentry inducibility in $182$ patient-derived
left-atrial meshes ($100$ Roney~\cite{roney2022}, $82$ University of
Washington~\cite{uwboyle2025}) and up to $100{,}000$ synthetic excitable
networks, together with a pre-registered clinical endpoint of two-year
post-ablation arrhythmia recurrence in the $82$ University of Washington
patients. The endpoints, cross-validation scheme, baseline feature set,
effect-size threshold ($\dAUC\geq 0.05$), and the accepted null that SFI merely
re-encodes fibrosis were frozen before any label was inspected.
